\section*{Appendix: reproducibility and claims register}
\phantomsection\label{sec:repro}
\addcontentsline{toc}{section}{Appendix: reproducibility and claims register}

\paragraph*{Measured chain}
\bfs commit \hash{\commitBeamfsFull}; Yocto layer \texttt{\commitYocto};
Linux 7.3-rc5, commit \texttt{72d3fcf802c45d00b300f25b848a93c3a2bd7c7e};
root images SHA-256: x86-64 \hash{\imageXSha}, aarch64 \hash{\imageSha};
kernels SHA-256: x86-64 \texttt{bzImage} \hash{\kernelXSha}, aarch64
\texttt{Image} \hash{\kernelASha}; \bfs built in on both
(\texttt{CONFIG\_BEAMFS\_FS=y}), \texttt{CONFIG\_BEAMFS\_DEBUG\_TREE=y}
on x86-64 only; \srcTrees{} compiled source trees identical to the
sources the manifests list; on x86-01, mkfs.beamfs \hash{\xfsXMkfsSha}
and fsck.beamfs \hash{\xfsXFsckSha}; emufi 0.8.1,
\texttt{srcversion 5DCD6301EB353B0FADD6FE2}; bench commits
\texttt{\commitBenchTwo}, \texttt{\commitBenchThree},
\texttt{\commitBenchFour} for runs 2, 3, 4.

\paragraph*{Signed manifests}
Each run manifest is signed with the author's key
\texttt{319A 8EAA 89C7 538A A955 0E8B C35E E212 519E 4857} (EdDSA) and
verified before extraction.
\begin{itemize}
\item run 2, \texttt{\manifestNameTwo}: \hash{\manifestShaTwo}
\item run 3, \texttt{\manifestNameThree}: \hash{\manifestShaThree}
\item run 4, \texttt{\manifestNameFour}: \hash{\manifestShaFour}
\end{itemize}

\paragraph*{xfstests records}
\begin{itemize}
\item x86-64 sweep, harness 2.5.1, \texttt{\xfsXtag}: \texttt{meta.txt}
  \hash{\xfsXmetaSha}, \texttt{verdicts.txt} \hash{\xfsXverdictsSha},
  check's output for each test, and the kernel log of each test in the
  harness's trace
\item aarch64 sweep, harness 2.3.60, \texttt{\sweepAfullName}:
  \hash{\sweepAfullSha}, and check's output for \xfsAcheckOutputs{} of
  its tests (\texttt{data/raw/xfstests-evidence})
\item x86-64 generic/476 alone, harness 2.3.60, \texttt{\sweepXsoloName}: \hash{\sweepXsoloSha}
\item what the harness read in the aarch64 volumes it froze when it
  stopped generic/476 (\texttt{data/raw/xfstests-476})
\item aarch64 generic/476 alone, \texttt{\sweepAsoloName}: \hash{\sweepAsoloSha}
\item superseded, not counted: x86-64 sweep, harness 2.3.60,
  \texttt{\sweepXfullName} (\hash{\sweepXfullSha}), saved as 734 passes,
  its per-test output overwritten; x86-64 sweep, harness 2.5.0,
  \texttt{\xfsSupTag} (\texttt{verdicts.txt} \hash{\xfsSupVerdictsSha}),
  stopped after \xfsSupDone{} tests (\texttt{data/raw/xfstests-superseded})
\end{itemize}

\paragraph*{Snapshot and scripts}
The whole workspace, \texttt{data/raw} included, is deposited with
this report, together with what the measurements ran on and what ran
them: the two root images and kernels above; for each architecture,
the kernel configuration, the image and license manifests, the Yocto
configuration with the commit of every layer, and the libvirt
definition of every node with its QEMU version; and the source trees,
at the commits used, of the Yocto layer, which carries beamfs,
\texttt{mkfs.beamfs}, \texttt{fsck.beamfs} and emufi~0.8.1 as built,
of the xfstests harness, of the bench, beamfs-bench, at the three
commits it ran from, and of the Gentoo ebuilds that install the
harness and the bench. The run
directories, manifests, bench database and verdict files were
copied once into \texttt{data/raw} (\snapshotFiles{} files) and frozen by
\texttt{data/raw/SHA256SUMS}, itself of SHA-256
\hash{\snapshotSha}. Three scripts turn them into the paper:
\texttt{scripts/extract.py} (\hash{\extractSha}) reads the raw records,
cross-checks them and writes \texttt{data/derived}; it writes nothing if
a check fails. \texttt{scripts/tables.py} (\hash{\tablesSha}) writes
every number and table the text quotes. \texttt{scripts/figures.jl}
(\hash{\figuresSha}) draws the figures. The commands, from the workspace
root:
\begin{lstlisting}
cd data/raw && sha256sum -c SHA256SUMS && cd ../..
for m in data/raw/runs/manifest-*.json; do gpg --verify "$m.asc" "$m"; done
python3 -I scripts/extract.py
python3 -I scripts/tables.py
julia scripts/figures.jl
latexmk -lualatex paper.tex
\end{lstlisting}

\paragraph*{Claims register}
Each claim, the record it comes from, and the check that guards it.
\begin{center}
\footnotesize
\begin{tabularx}{\textwidth}{@{}>{\raggedright\arraybackslash}p{4.6cm}
  >{\raggedright\arraybackslash}p{4.6cm} >{\raggedright\arraybackslash}X@{}}
\toprule
claim & source & check \\
\midrule
same \bfs commit, layer and image in runs 2 to 4 & three signed manifests &
  GPG signature; equality across manifests (\texttt{tables.py}) \\
same \bfs sources, images and kernels on both architectures & seals,
  deploy directories, compiled trees (\texttt{identity.txt}) & image
  equal to seal and to the manifests; trees identical
  (\texttt{extract.py}) \\
xfstests: \xfsXPass{} of \xfsXRan{} run on x86-64, \xfsAPass{} of
  \xfsARan{} on aarch64, declined tests apart & check's output for each
  test & every verdict read from it; on x86-64 equal to the harness's
  (\texttt{extract.py}) \\
most declined tests for features \bfs lacks & reason check gave per
  test & every reason classified; feature share above half
  (\texttt{tables.py}) \\
generic/476: passes in the x86-64 sweep; stopped by the budget in the
  aarch64 sweep, passes alone there & harness results and records &
  verdicts; seconds within the budget (\texttt{extract.py},
  \texttt{tables.py}) \\
budget of the x86-64 sweep, 14\,400\,s; superseded sweep stopped within
  1\,900\,s & \texttt{meta.txt} and \texttt{verdicts.txt} of each record &
  \texttt{budget=} lines; shell's kill report (\texttt{extract.py},
  \texttt{tables.py}) \\
frozen aarch64 volumes of 476: pointers beyond the device & harness's
  reading of the images & phrase present (\texttt{extract.py}) \\
no warning or oops in the x86-64 sweep, no lost pointer & taint, uptime
  and recoveries in \texttt{meta.txt}; kernel log of each test & taint
  equal, warning and oops bits clear, no restart; no such line in any
  log (\texttt{extract.py}, \texttt{tables.py}) \\
btrfs refused on a checksum mismatch &
  \nolinkurl{forensics-beamfs-compute01/dmesg.log} of each run &
  \texttt{csum failed} lines present; none in the bench records
  (\texttt{extract.py}) \\
flips and outcome of every filesystem attack & \texttt{multifs-all-records.txt}
  of each run & identical to the multifs directory copy; call, flip,
  read status, integrity and bit counts equal to the bench database \\
outcome of every cluster attack & \texttt{cluster-records.txt} &
  outcome recomputed from the record equals the bench verdict \\
flips of an attack are the last FLIP\_DELTA ring entries & flip ring in
  each record & ring length equals the running sum of FLIP\_DELTA, every
  attack \\
two hooks per bio; pairs within \dtPairMaxMicro\,$\mu$s, others
  $\geq$\dtOtherMinMilli\,ms & flip ring & no interval between; every
  flip at $10^{6}$\,ppm attributed \\
all \kernLines{} kernel correction lines explained & kernel log slice in
  each record & matching on block and codeword, no line left \\
at most \kernMaxPerDecode{} symbols per decode, \unionMax{} per codeword
  per attack & kernel log; reconstructed flips & bound $\leq 8$
  (\texttt{extract.py}) \\
indirect block hit at slots \indSlots{}, none in use & reconstructed flips
  & slot $<$ 57 counted (\texttt{tables.py}) \\
run 3 pointer flip, slot 33, bit 54 & \nolinkurl{forensics-beamfs-master/dmesg.log}
  of run 3 & line present four times (\texttt{check-paper.sh}) \\
\texttt{incompat} \featIncompat{}, \texttt{ro\_compat} 0 & mount line in
  each kernel log slice & equal in all six \bfs attacks (\texttt{tables.py}) \\
medium unaltered & verification lines & twelve files intact after every
  attack, all filesystems; all \bfs flips on reads \\
module of \moduleLines{} lines & \texttt{git show} at \texttt{\commitBeamfs}
  & recomputed at each build (\texttt{tables.py}) \\
\texttt{bio\_advance} leaves the sector of a complete bio & kernel source
  at \texttt{72d3fcf8} & function body checked (\texttt{check-paper.sh}) \\
statements on Memon et al. & \cite{memon2025tr,memon2025arxiv}, read
  2026-10-08 & quoted in the text \\
\bottomrule
\end{tabularx}
\end{center}
